\makereporttitle{2024}{octombrie}

\makereportpreamble

În această lună am început să merg în mod regulat la cursuri de master de la Facultatea de Fizică și de la Facultatea de Matematică și Informatică, pentru a aprofunda anumite noțiuni pe care le-am identificat ca fiind de interes pentru problemele de cercetare pe care lucrez.

Am participat la cursul de ``\textbf{AI in Theoretical Physics}'' organizat de domnul profesor Ovidiu Călin (\emph{visiting scholar}, întors în țară printr-o bursă Fulbright) în cadrul masterului de \emph{Theoretical and Computational Physics}, unde au fost prezentate (într-o manieră interactivă și inovativă) noțiuni fundamentale referitoare la metodele moderne de învățare automată.

Am mers, de asemenea, la câteva cursuri de la masterele de \emph{Advanced Studies in Mathematics} și \emph{Probabilități și Statistică în Finanțe și Științe}: cursul de ``\textbf{Procese stocastice}'', ținut de dl.~prof.~Ionel Popescu; cursul de ``\textbf{Metode numerice pentru sisteme non-liniare și optimizare}'', ținut de dl.~prof.~Liviu Marin; cursul de ``\textbf{Tehnici de programare convexă}'', ținut de dl.~prof.~Cristian Niculescu; precum și cursul de ``\textbf{Riemannian Geometry}'', ținut de dl.~prof.~Sergiu Moroianu. De câteva ori am mers și la cursul de licență ``\textbf{Calcul Numeric în Informatică}'', ținut de conferențiarii Andrei Pătrașcu, Cristian Rusu și Paul Irofti.

În perioada 7-11 octombrie a avut loc online a o suta sesiune a ``\href{https://edpemaplicada.uff.br/wp-content/uploads/sites/308/2024/10/book_of_abstracts_SPDEAM24-1.pdf}{Seminar in PDEs and Applied Mathematics}'', unde au fost prezentate diferite topici referitoare la starea tehnicii în domeniul ecuațiilor diferențiale și a matematici.

Pe data de 15 octombrie 2024, am participat la un workshop organizat de Institutul de Statistică Matematică și Matematică Aplicată ``Gheorghe Mihoc -- Caius Iacob'' al Academiei Române, cu ocazia aniversării a 60 de ani de activitate a Centrului de Statistică Matematică. Mai mulți cercetători au luat cuvântul pentru a prezenta zonele lor de activitate și stadiul curent al cercetării matematice din România.

Pe data de 23 octombrie 2024, am participat la o conferință la Academia Română, dedicată lui David Emmanuel, considerat de unii drept fondatorul școlii românești de matematică.

În paralel, am continuat să învăț cum să implementez algoritmi numerici folosind limbajele Python, Julia și MATLAB.

\vspace{0.5em}

\textbf{Observațiile conducătorului științific}:

\vspace{3em}

\makereportpostamble{30 octombrie 2024}
